%==================================appendix.tex=============================
\newpage
\begin{spacing}{1.5}
%
\chapter*{APPENDIX}
%\addcontentsline{toc}{chapter}{\numberline{}Appendix}%
\addcontentsline{toc}{chapter}{Appendix}%

\thispagestyle{empty}
%%
\noindent \textbf{CODE SNIPPETS}\\

%% ===================================================================
%% SECTION A: BACKEND SERVER (Node.js)
%% ===================================================================

\noindent \textbf{A. Backend Server (Node.js)}\\

\noindent \textbf{A.1 JWT Authentication Middleware}\\
Description: Protects API routes by verifying JSON Web Tokens from incoming requests. The decoded user payload (userId, role, username) is attached to the request for downstream use.
\begin{verbatim}
const authMiddleware = (req, res, next) => {
  const authHeader = req.headers.authorization;
  if (!authHeader || !authHeader.startsWith("Bearer "))
    return res.status(401).json({ message: "No token" });

  const token = authHeader.split(" ")[1];
  try {
    req.user = jwt.verify(token, process.env.JWT_SECRET);
    next();
  } catch (err) {
    return res.status(401).json({ message: "Invalid token" });
  }
};
\end{verbatim}

\noindent \textbf{A.2 User Intent Classification (Gemini AI)}\\
Description: Classifies the user's transcribed speech into one of seven assistance modules using a Google Gemini prompt. The returned module name drives the Pi client's next action.
\begin{verbatim}
async function identifyUserIntent(transcribedText) {
  const intentPrompt = `You are an intent classifier.
Available modules:
1. "Product Identification Module"
2. "Medical Compatibility Module"
3. "Currency Recognition Module"
4. "Volunteer Video Call Module"
5. "AI Assistance Module"
6. "Scene Description Module"
7. "Price Comparison Module"
Respond with ONLY the exact module name.
User said: "${transcribedText}"`;

  const resp = await gemini.models.generateContent({
    model: "gemini-2.5-flash", contents: intentPrompt
  });
  return (resp.text || "AI Assistance Module").trim();
}
\end{verbatim}

\noindent \textbf{A.3 Intent-to-Action Mapping (/pi\_intent Route)}\\
Description: Maps the detected intent module to an action command that the Raspberry Pi client executes. For AI conversation intents, the response is generated inline via Gemini.
\begin{verbatim}
switch (detectedModule) {
  case "Product Identification Module":
    actionCommand = "CAPTURE_PRODUCT_IMAGE"; break;
  case "Medical Compatibility Module":
    actionCommand = "CAPTURE_MEDICAL_IMAGE"; break;
  case "Currency Recognition Module":
    actionCommand = "CAPTURE_CURRENCY_IMAGE"; break;
  case "Volunteer Video Call Module":
    actionCommand = "INITIATE_VIDEO_CALL"; break;
  case "Scene Description Module":
    actionCommand = "CAPTURE_ENVIRONMENT"; break;
  default:
    actionCommand = "AI_CONVERSATION"; break;
}
\end{verbatim}
\thispagestyle{empty}
\noindent \textbf{A.4 Product Identification using GPT Vision}\\
Description: Sends a Base64-encoded product image to the OpenAI GPT-4 Vision API, which returns only the brand and product name.
\begin{verbatim}
async function identifyProductWithGPT(imageBuffer, mimeType) {
  const base64 = imageBuffer.toString("base64");
  const resp = await client.responses.create({
    model: "gpt-4.1-mini",
    input: [{
      role: "user",
      content: [
        { type: "input_text",
          text: "Reply with ONLY the product name." },
        { type: "input_image",
          image_url: `data:${mimeType};base64,${base64}` }
      ]
    }]
  });
  return (resp.output_text || "").trim();
}
\end{verbatim}

\noindent \textbf{A.5 Medical Compatibility Check (Core Pipeline)}\\
Description: Identifies a product from its image, fetches ingredient data via SerpAPI, and queries Gemini for medical consumability advice based on the user's stored health profile.
\begin{verbatim}
// Step 1: Identify product from image
const productName = await identifyProductWithGPT(
  req.file.buffer, req.file.mimetype);

// Step 2: Search ingredients via SerpAPI
const serpData = await searchSerpApi(
  `${productName} ingredients and allergens`);

// Step 3: Get consumability advice from Gemini
const aiAdvice = await getConsumabilityAdviceWithGemini({
  productName, serpData,
  medicalProfile: blindProfile
});

return res.status(200).json({
  product_name: productName, advice: aiAdvice
});
\end{verbatim}

\noindent \textbf{A.6 Volunteer Matching (Atomic Lock)}\\
Description: When a blind user requests a video call, the server atomically finds and locks an available volunteer, generates a deterministic LiveKit room ID, and returns connection credentials.
\begin{verbatim}
const volunteer = await VolunteerProfile.findOneAndUpdate(
  { isAvailable: true, consentGiven: true },
  { isAvailable: false },
  { new: true }
).populate("user", "username email _id");

const roomID = generateRoomId(blindId, volunteerId);
res.json({ success: true, roomID, volunteer });
\end{verbatim}
\thispagestyle{empty}
%% ===================================================================
%% SECTION B: RASPBERRY PI CLIENT (Python)
%% ===================================================================

\noindent \textbf{B. Raspberry Pi Client (Python)}\\

\noindent \textbf{B.1 Hardware Configuration}\\
Description: Defines the hardware interface constants for the INMP441 I$^2$S microphone, TTP223 touch sensor, and audio recording parameters on the Raspberry Pi Zero 2 W.
\begin{verbatim}
TOUCH_SENSOR_PIN = 17    # GPIO17 - TTP223 OUT
AUDIO_DEVICE     = "hw:0,0"
SAMPLE_RATE      = 16000
CHANNELS         = 2
RECORD_FORMAT    = "S32_LE"  # 32-bit signed LE
MIN_RECORD_DURATION = 0.5   # seconds
MAX_RECORD_DURATION = 30.0  # seconds
SILENCE_THRESHOLD   = 2.0   # seconds
\end{verbatim}

\noindent \textbf{B.2 Touch-Triggered Audio Recording (Key Logic)}\\
Description: Records audio via \texttt{arecord} while the touch sensor is held. The recording loop polls the GPIO pin at 50ms intervals and stops after the sensor is released for the silence threshold duration.
\begin{verbatim}
process = subprocess.Popen(
    ["arecord", "-D", AUDIO_DEVICE, "-f", RECORD_FORMAT,
     "-r", str(SAMPLE_RATE), "-c", str(CHANNELS),
     "-t", "wav", output_path])

while True:
    elapsed = time.time() - start_time
    if elapsed >= MAX_RECORD_DURATION:
        break

    touch_active = GPIO.input(TOUCH_SENSOR_PIN) == GPIO.HIGH
    if touch_active:
        touch_released_time = None
    else:
        if touch_released_time is None:
            touch_released_time = time.time()
        elif (time.time() - touch_released_time
                >= SILENCE_THRESHOLD):
            break

    time.sleep(0.05)  # 50ms polling

process.terminate()
\end{verbatim}

\noindent \textbf{B.3 Speech-to-Text via FastAPI}\\
Description: Sends the recorded WAV file to the FastAPI speech-to-text service and returns the transcribed text.
\begin{verbatim}
def speech_to_text(audio_path: str) -> str:
    with open(audio_path, 'rb') as audio_file:
        files = {'file': ('recording.wav',
                          audio_file, 'audio/wav')}
        response = requests.post(
            f"{FASTAPI_URL}/stt-upload",
            files=files, timeout=60)

    result = response.json()
    return result.get('text', '').strip() or None
\end{verbatim}
\thispagestyle{empty}
\noindent \textbf{B.4 Main Interaction Workflow (Dispatch Logic)}\\
Description: After recording and transcribing the user's voice, the client sends the text to the server for intent detection. Based on the returned action command, it dispatches to the appropriate handler --- capturing images, initiating video calls, or playing AI conversation responses.
\begin{verbatim}
# Record audio and transcribe
success = record_with_touch_trigger(recording_path)
transcribed_text = speech_to_text(recording_path)
intent_data = send_text_to_intent(transcribed_text)

action = intent_data.get('action_command', '')

if action == "CAPTURE_PRODUCT_IMAGE":
    capture_image(image_path)
    audio = send_image_to_product_identification(image_path)

elif action == "CAPTURE_MEDICAL_IMAGE":
    capture_image(image_path)
    audio = send_image_to_medical_check(image_path)

elif action == "CAPTURE_CURRENCY_IMAGE":
    capture_image(image_path)
    audio = send_image_to_currency_recognition(image_path)

elif action == "INITIATE_VIDEO_CALL":
    asyncio.run(initiate_video_call())

elif action == "AI_CONVERSATION":
    ai_text = intent_data.get('ai_response')
    audio = requests.post(
        f"{FASTAPI_URL}/tts",
        json={"text": ai_text}).content

if audio:
    play_audio_pulseaudio(audio)
\end{verbatim}

\noindent \textbf{B.5 GPIO Setup and Main Loop}\\
Description: Initializes the TTP223 touch sensor on GPIO17 and runs the main event loop, which waits for touch activation and processes each interaction.
\begin{verbatim}
def setup_gpio():
    GPIO.setmode(GPIO.BCM)
    GPIO.setup(TOUCH_SENSOR_PIN, GPIO.IN)

def main():
    if not login():
        return
    setup_gpio()
    try:
        while True:
            if not wait_for_touch():
                break
            process_single_interaction()
            time.sleep(0.5)
    except KeyboardInterrupt:
        pass
    finally:
        GPIO.cleanup()

if __name__ == "__main__":
    main()
\end{verbatim}

%% ===================================================================
%% SECTION C: FASTAPI ML SERVER (Python)
%% ===================================================================

\noindent \textbf{C. FastAPI ML Server (Python)}\\

\noindent \textbf{C.1 Speech-to-Text using OpenAI Whisper}\\
Description: Receives an uploaded audio file, saves it temporarily, and transcribes it using the Whisper \texttt{turbo} model with forced English language detection. The transcribed text is returned as JSON.
\begin{verbatim}
whisper_model = whisper.load_model("turbo")

@app.post("/stt-upload")
async def stt_upload(file: UploadFile = File(...)):
    with tempfile.NamedTemporaryFile(
        delete=False, suffix=suffix) as tmp:
        contents = await file.read()
        tmp.write(contents)
        temp_path = tmp.name

    result = whisper_model.transcribe(
        temp_path, language="en")
    text = result.get("text", "")
    os.remove(temp_path)
    return {"text": text}
\end{verbatim}

\noindent \textbf{C.2 Text-to-Speech using Kokoro TTS}\\
Description: Converts input text to natural-sounding speech audio using the Kokoro TTS pipeline. The generated audio chunks are concatenated and returned as a streaming WAV response.
\begin{verbatim}
pipeline = KPipeline(lang_code="a")

@app.post("/tts")
def tts(text: str = Body(..., embed=True)):
    generator = pipeline(text, voice="af_heart")

    chunks = []
    for _, _, audio in generator:
        chunks.append(audio)

    full_audio = np.concatenate(chunks)
    buffer = io.BytesIO()
    sf.write(buffer, full_audio, 24000, format="WAV")
    buffer.seek(0)
    return StreamingResponse(
        buffer, media_type="audio/wav")
\end{verbatim}
\thispagestyle{empty}
\noindent \textbf{C.3 Currency Recognition (ConvNeXt + Prototypical Network)}\\
Description: Uses a ConvNeXt-Tiny backbone with a learned embedding projection head for few-shot currency recognition. The model compares the input image embedding against stored class prototypes using cosine similarity with a learned temperature scale.
\begin{verbatim}
class ConvNeXtIncremental(nn.Module):
    def __init__(self, embedding_dim=256):
        super().__init__()
        backbone = torchvision.models.convnext_tiny()
        self.features = backbone.features
        self.pool = nn.AdaptiveAvgPool2d(1)
        in_features = backbone.classifier[2].in_features
        self.proj = nn.Sequential(
            nn.Linear(in_features, 512),
            nn.LayerNorm(512), nn.ReLU(),
            nn.Linear(512, embedding_dim))

    def forward(self, x):
        x = self.features(x)
        x = self.pool(x).flatten(1)
        return F.normalize(self.proj(x), p=2, dim=1)

@app.post("/currency-recognition")
async def currencyRecognition(file: UploadFile):
    img = Image.open(io.BytesIO(await file.read()))
    x = currency_transform(img).unsqueeze(0).to(DEVICE)

    with torch.no_grad():
        emb = currency_model(x)
        sims = torch.mm(emb, currency_prototypes.t())
        probs = F.softmax(sims * currency_scale, dim=1)

    return {
        "prediction": currency_classes[
            probs.argmax(1).item()],
        "confidence": float(probs.max().item())
    }
\end{verbatim}

\noindent \textbf{C.4 YOLO-Based Object Detection and Cropping}\\
Description: Runs YOLOv8 inference on uploaded images to detect objects. Each detection is cropped with configurable padding and saved as an individual image for downstream processing.
\begin{verbatim}
yolo_model = YOLO(MODEL_PATH)

@app.post("/detect-crop")
async def detect_crop(file: UploadFile = File(...)):
    data = await file.read()
    image = cv2.imdecode(
        np.frombuffer(data, np.uint8),
        cv2.IMREAD_COLOR)

    results = yolo_model.predict(
        source=image, conf=0.25)

    for result in results:
        for box in result.boxes:
            x1, y1, x2, y2 = map(int, box.xyxy[0])
            # Apply padding and crop
            cropped = image[y1p:y2p, x1p:x2p]
            cv2.imwrite(str(crop_path), cropped)

    return {"ok": True, "crops": crop_count}
\end{verbatim}

%%
%%%%%%%%%%%%%%%%%%%%%%%%%%%%%%%%%%%%%
%%


\end{spacing}
